\documentclass[12pt]{article}
\usepackage[T1]{fontenc}
\usepackage[utf8]{inputenc}
\usepackage{lmodern}
\usepackage{textcomp}
\usepackage{enumitem}
\usepackage{geometry}
\geometry{margin=1in}
\begin{document}
\begin{center}
Answer Key - Political Science - Graduate Level Assessment 16 \
\small (Answers are ordered exactly as in the question paper.)
\end{center}

\section*{Multiple Choice}
\begin{enumerate}[label=\arabic*.]
\item \textbf{Answer:} A
\\ \textit{Options:} A) Realism; B) Idealism; C) Liberalism; D) None of the above
\\ \textit{Note:} Realism in political science emphasizes the idea that states should prioritize their national interests and security over moral considerations, arguing that the use of force abroad for humanitarian or democratic purposes can lead to unintended consequences and destabilization.
\item \textbf{Answer:} D
\\ \textit{Options:} A) James Madison; B) Abraham Lincoln; C) Woodrow Wilson; D) Thomas Jefferson
\\ \textit{Note:} This statement reflects Thomas Jefferson's foreign policy vision articulated in his inaugural address, emphasizing the importance of maintaining peaceful and commercial relationships with other nations while avoiding complicated alliances that could lead to conflict.
\item \textbf{Answer:} C
\\ \textit{Options:} A) the importance of military superiority.; B) how the importance of oil is often overexaggerated.; C) the increasing importance of economic instruments of foreign policy.; D) the need to drill for domestic sources.
\\ \textit{Note:} The correct answer emphasizes that the reliance on oil has driven the U.S. to prioritize economic interests, illustrating how economic instruments, such as sanctions and trade agreements, have become central to its foreign policy strategies in securing energy resources and maintaining global market stability.
\item \textbf{Answer:} A
\\ \textit{Options:} A) A new democratic internationalism led by the United States; B) A new balance of power between the US and China; C) A new global economic framework; D) A new era of globalization
\\ \textit{Note:} The term 'New World Order' refers to the post-Cold War vision of a more integrated and cooperative global system characterized by democracy and multilateralism, primarily promoted by the United States as it sought to expand its influence and establish a framework for international relations.
\item \textbf{Answer:} A
\\ \textit{Options:} A) It allowed the US to intensify its involvement in Vietnam; B) It illustrated the influence of public opinion on US foreign policy; C) It enhanced Congressional control over the Vietnam War; D) It curtailed US involvement in Vietnam
\\ \textit{Note:} The Gulf of Tonkin resolution was significant because it granted President Lyndon B. Johnson broad authority to escalate military action in Vietnam without a formal declaration of war, effectively leading to increased U.S. involvement in the conflict.
\end{enumerate}

\section*{True / False}
\begin{enumerate}[label=\arabic*.]
\setcounter{enumi}{5}
\item \textbf{Answer:} False
\\ \textit{Options:} A) True; B) False
\\ \textit{Note:} The argument for American unilateralism is based on the belief that the U.S. can act independently and more effectively to protect its interests, rather than on the successes of multilateralism, which often emphasizes collaboration and compromise among multiple nations.
\item \textbf{Answer:} False
\\ \textit{Options:} A) True; B) False
\\ \textit{Note:} The grand strategy of 'primacy' emphasizes the United States maintaining a dominant position over other global powers, rather than merely achieving parity.
\item \textbf{Answer:} False
\\ \textit{Options:} A) True; B) False
\\ \textit{Note:} The director of the CIA primarily focuses on intelligence gathering and analysis rather than serving as the president's principal civilian adviser on military issues, which is typically the role of the Secretary of Defense.
\item \textbf{Answer:} True
\\ \textit{Options:} A) True; B) False
\\ \textit{Note:} Bureaucratic politics highlights the necessity of understanding the roles and interests of various political actors and institutions in decision-making processes, particularly in matters of national security, such as identifying and responding to nuclear threats and ensuring effective control over nuclear arsenals.
\item \textbf{Answer:} True
\\ \textit{Options:} A) True; B) False
\\ \textit{Note:} Post-1865, the United States shifted its foreign policy focus from acquiring new territories to promoting economic interests and trade relationships, exemplified by policies such as the Open Door Policy and the emphasis on securing access to markets in Asia and Latin America.
\end{enumerate}

\section*{Long Form Response}
\begin{enumerate}[label=\arabic*.]
\setcounter{enumi}{10}
\item \textbf{Answer:} NSC 68 marked a significant transformation in U.S. Cold War foreign policy by advocating for a robust military containment strategy aimed at countering Soviet expansionism. This document underscored the necessity for an increase in defense spending and a more aggressive posture towards communism, reflecting a shift from previous diplomatic approaches to one that prioritized military readiness and deterrence. The implications of NSC 68 were profound, as it not only set the stage for the militarization of U.S. foreign policy but also solidified a worldview that framed international relations in stark terms of ideological conflict. By positioning the U.S. as a bulwark against the perceived threat of communism, NSC 68 influenced subsequent actions such as the Korean War and the establishment of NATO, thereby reshaping alliances and enmities in the global arena.
\\ \textit{Note:} correct because it identifies NSC 68 as a critical turning point in U.S. Cold War policy that emphasized military containment and preparedness against Soviet threats, highlighting a shift from diplomatic strategies to a focus on military strength and its significant impact on international relations.
\item \textbf{Answer:} The Western response to the Arab Spring presented significant challenges to liberalism, as it revealed inconsistencies in the promotion of democratic values and human rights. Initially, the West supported uprisings against authoritarian regimes; however, when stability was threatened by the rise of Islamist groups or civil unrest, there was a tendency to revert to supporting autocratic governments. This selective engagement undermines the credibility of liberalism by suggesting that the West prioritizes strategic interests over genuine democratic governance. Moreover, the reaction to the Arab Spring has strained international relations, as it has fostered skepticism towards Western intentions in the Middle East, complicating future diplomatic efforts. Ultimately, the complexities of these responses highlight the fragility of liberal ideals in the face of geopolitical realities.
\\ \textit{Note:} correct because it demonstrates how the West's inconsistent support for democratic movements during the Arab Spring-favoring stability over democratic values when faced with Islamist groups or unrest-undermined the foundational principles of liberalism and complicated the broader framework of international relations and democratic governance.
\end{enumerate}

\vfill
\noindent\textit{End of Answer Key}
\end{document}